\section{Cohomology of projective space}
\label{section-cohomology-projective-space}

\noindent
In this section we compute the cohomology of the twists of the
structure sheaf on $\mathbf{P}^n_S$ over a scheme $S$.
Recall that $\mathbf{P}^n_S$ was defined as the fibre product
$
\mathbf{P}^n_S = S \times_{\Spec(\mathbf{Z})} \mathbf{P}^n_{\mathbf{Z}}
$
in Constructions, Definition \ref{constructions-definition-projective-space}.
It was shown to be equal to
$$
\mathbf{P}^n_S = \underline{\text{Proj}}_S(\mathcal{O}_S[T_0, \ldots, T_n])
$$
in Constructions, Lemma \ref{constructions-lemma-projective-space-bundle}.
In particular, projective space is a particular case of a projective bundle.
If $S = \Spec(R)$ is affine then we have
$$
\mathbf{P}^n_S = \mathbf{P}^n_R = \text{Proj}(R[T_0, \ldots, T_n]).
$$
All these identifications are compatible and compatible with the constructions
of the twisted structure sheaves $\mathcal{O}_{\mathbf{P}^n_S}(d)$.

\medskip\noindent
Before we state the result we need some notation.
Let $R$ be a ring.
Recall that $R[T_0, \ldots, T_n]$ is a graded
$R$-algebra where each $T_i$ is homogeneous of degree $1$.
Denote $(R[T_0, \ldots, T_n])_d$ the degree $d$ summand.
It is a finite free $R$-module of rank $\binom{n + d}{d}$
when $d \geq 0$ and zero else.
It has a basis consisting of monomials $T_0^{e_0} \ldots T_n^{e_n}$
with $\sum e_i = d$. We will also use the following notation:
$R[\frac{1}{T_0}, \ldots, \frac{1}{T_n}]$ denotes the $\mathbf{Z}$-graded
ring with $\frac{1}{T_i}$ in degree $-1$. In particular the
$\mathbf{Z}$-graded $R[\frac{1}{T_0}, \ldots, \frac{1}{T_n}]$ module
$$
\frac{1}{T_0 \ldots T_n} R[\frac{1}{T_0}, \ldots, \frac{1}{T_n}]
$$
which shows up in the statement below is zero in degrees
$\geq -n$, is free on the generator $\frac{1}{T_0 \ldots T_n}$
in degree $-n - 1$ and is free of rank $(-1)^n\binom{n + d}{d}$ for
$d \leq -n - 1$.

\begin{lemma}
\label{lemma-cohomology-projective-space-over-ring}
\begin{reference}
\cite[III Proposition 2.1.12]{EGA}
\end{reference}
Let $R$ be a ring.
Let $n \geq 0$ be an integer.
We have
$$
H^q(\mathbf{P}^n, \mathcal{O}_{\mathbf{P}^n_R}(d)) =
\left\{
\begin{matrix}
(R[T_0, \ldots, T_n])_d & \text{if} & q = 0,\ d \geq 0 \\
\left(\frac{1}{T_0 \ldots T_n} R[\frac{1}{T_0}, \ldots, \frac{1}{T_n}]\right)_d
& \text{if} & q = n,\ d < 0 \\
0 & \text{if} & q, n, d\text{ not as above}
\end{matrix}
\right.
$$
as $R$-modules.
\end{lemma}

\begin{proof}
We will use the standard affine open covering
$$
\mathcal{U} : \mathbf{P}^n_R = \bigcup\nolimits_{i = 0}^n D_{+}(T_i)
$$
to compute the cohomology using the {\v C}ech complex.
This is permissible by Lemma \ref{lemma-cech-cohomology-quasi-coherent}
since any intersection of finitely many affine $D_{+}(T_i)$ is also a
standard affine open (see
Constructions, Section \ref{constructions-section-proj}).
In fact, we can use the alternating or ordered {\v C}ech complex according to
Cohomology, Lemmas \ref{cohomology-lemma-ordered-alternating} and
\ref{cohomology-lemma-alternating-usual}.

\medskip\noindent
The ordering we will use on $\{0, \ldots, n\}$ is the usual one.
Hence the complex we are looking at has terms
$$
\check{\mathcal{C}}_{ord}^p(\mathcal{U}, \mathcal{O}_{\mathbf{P}_R}(d))
=
\bigoplus\nolimits_{i_0 < \ldots < i_p}
(R[T_0, \ldots, T_n, \frac{1}{T_{i_0} \ldots T_{i_p}}])_d
$$
Moreover, the maps are given by the usual formula
$$
d(s)_{i_0 \ldots i_{p + 1}} =
\sum\nolimits_{j = 0}^{p + 1} (-1)^j s_{i_0 \ldots \hat i_j \ldots i_{p + 1}}
$$
see Cohomology, Section \ref{cohomology-section-alternating-cech}.
Note that each term of this complex has a natural
$\mathbf{Z}^{n + 1}$-grading. Namely, we get this by declaring a monomial
$T_0^{e_0} \ldots T_n^{e_n}$ to be homogeneous with weight
$(e_0, \ldots, e_n) \in \mathbf{Z}^{n + 1}$. It is clear that the differential
given above respects the grading. In a formula we have
$$
\check{\mathcal{C}}_{ord}^\bullet(\mathcal{U}, \mathcal{O}_{\mathbf{P}_R}(d))
=
\bigoplus\nolimits_{\vec{e} \in \mathbf{Z}^{n + 1}}
\check{\mathcal{C}}^\bullet(\vec{e})
$$
where not all summands on the right hand side occur (see below).
Hence in order to compute the cohomology
modules of the complex it suffices to compute the cohomology of the graded
pieces and take the direct sum at the end.

\medskip\noindent
Fix $\vec{e} = (e_0, \ldots, e_n) \in \mathbf{Z}^{n + 1}$. In order for this
weight to occur in the complex above we need to assume
$e_0 + \ldots + e_n = d$ (if not then it occurs for a different twist of
the structure sheaf of course). Assuming this, set
$$
NEG(\vec{e}) = \{i \in \{0, \ldots, n\} \mid e_i < 0\}.
$$
With this notation the weight $\vec{e}$ summand
$\check{\mathcal{C}}^\bullet(\vec{e})$ of the {\v C}ech complex above has
the following terms
$$
\check{\mathcal{C}}^p(\vec{e})
=
\bigoplus\nolimits_{i_0 < \ldots < i_p,
\ NEG(\vec{e}) \subset \{i_0, \ldots, i_p\}}
R \cdot T_0^{e_0} \ldots T_n^{e_n}
$$
In other words, the terms corresponding to $i_0 < \ldots < i_p$ such
that $NEG(\vec{e})$ is not contained in $\{i_0 \ldots i_p\}$ are zero.
The differential of the complex $\check{\mathcal{C}}^\bullet(\vec{e})$
is still given by the exact same formula as above.

\medskip\noindent
Suppose that $NEG(\vec{e}) = \{0, \ldots, n\}$, i.e., that all
exponents $e_i$ are negative.
In this case the complex $\check{\mathcal{C}}^\bullet(\vec{e})$ has
only one term, namely $\check{\mathcal{C}}^n(\vec{e}) =
R \cdot \frac{1}{T_0^{-e_0} \ldots T_n^{-e_n}}$. Hence in this
case
$$
H^q(\check{\mathcal{C}}^\bullet(\vec{e})) =
\left\{
\begin{matrix}
R \cdot \frac{1}{T_0^{-e_0} \ldots T_n^{-e_n}} & \text{if} & q = n \\
0 & \text{if} & \text{else}
\end{matrix}
\right.
$$
The direct sum of all of these terms clearly gives the value
$$
\left(\frac{1}{T_0 \ldots T_n} R[\frac{1}{T_0}, \ldots, \frac{1}{T_n}]\right)_d
$$
in degree $n$ as stated in the lemma. Moreover these terms do not contribute
to cohomology in other degrees (also in accordance with the statement of the
lemma).

\medskip\noindent
Assume $NEG(\vec{e}) = \emptyset$. In this case the complex
$\check{\mathcal{C}}^\bullet(\vec{e})$ has a summand $R$ corresponding
to all $i_0 < \ldots < i_p$.
Let us compare the complex $\check{\mathcal{C}}^\bullet(\vec{e})$
to another complex. Namely, consider the affine open covering
$$
\mathcal{V} : \Spec(R) = \bigcup\nolimits_{i \in \{0, \ldots, n\}} V_i
$$
where $V_i = \Spec(R)$ for all $i$. Consider the alternating
{\v C}ech complex
$$
\check{\mathcal{C}}_{ord}^\bullet(\mathcal{V}, \mathcal{O}_{\Spec(R)})
$$
By the same reasoning as above this computes the cohomology of the
structure sheaf on $\Spec(R)$. Hence we see that
$H^p(
\check{\mathcal{C}}_{ord}^\bullet(\mathcal{V}, \mathcal{O}_{\Spec(R)})
) = R$ if $p = 0$ and is $0$ whenever $p > 0$.
For these facts, see
Lemma \ref{lemma-cech-cohomology-quasi-coherent-trivial} and its proof.
Note that also
$\check{\mathcal{C}}_{ord}^\bullet(\mathcal{V}, \mathcal{O}_{\Spec(R)})$
has a summand $R$ for every $i_0 < \ldots < i_p$ and has exactly the same
differential as $\check{\mathcal{C}}^\bullet(\vec{e})$. In other words
these complexes are isomorphic complexes and hence have the same cohomology.
We conclude that
$$
H^q(\check{\mathcal{C}}^\bullet(\vec{e})) =
\left\{
\begin{matrix}
R \cdot T_0^{e_0} \ldots T_n^{e_n} & \text{if} & q = 0 \\
0 & \text{if} & \text{else}
\end{matrix}
\right.
$$
in the case that $NEG(\vec{e}) = \emptyset$.
The direct sum of all of these terms clearly gives the value
$$
(R[T_0, \ldots, T_n])_d
$$
in degree $0$ as stated in the lemma. Moreover these terms do not contribute
to cohomology in other degrees (also in accordance with the statement of the
lemma).

\medskip\noindent
To finish the proof of the lemma we have to show that the complexes
$\check{\mathcal{C}}^\bullet(\vec{e})$ are acyclic when
$NEG(\vec{e})$ is neither empty nor equal to $\{0, \ldots, n\}$.
Pick an index $i_{\text{fix}} \not \in NEG(\vec{e})$ (such an index exists).
Consider the map
$$
h :
\check{\mathcal{C}}^{p + 1}(\vec{e})
\to
\check{\mathcal{C}}^p(\vec{e})
$$
given by the rule that for $i_0 < \ldots < i_p$ we have
$$
h(s)_{i_0 \ldots i_p} =
\left\{
\begin{matrix}
0 & \text{if} & p \not \in \{0, \ldots, n - 1\} \\
0 & \text{if} & i_{\text{fix}} \in \{i_0, \ldots, i_p\} \\
s_{i_{\text{fix}} i_0 \ldots i_p} & \text{if} & i_{\text{fix}} < i_0 \\
(-1)^a s_{i_0 \ldots i_{a - 1} i_{\text{fix}} i_a \ldots i_p} &
\text{if} & i_{a - 1} < i_{\text{fix}} < i_a \\
(-1)^{p + 1} s_{i_0 \ldots i_p i_{\text{fix}}} &
\text{if} & i_p < i_{\text{fix}}
\end{matrix}
\right.
$$
Please compare with the proof of
Lemma \ref{lemma-cech-cohomology-quasi-coherent-trivial}.
This makes sense because we have
$$
NEG(\vec{e}) \subset \{i_0, \ldots, i_p\}
\Leftrightarrow
NEG(\vec{e}) \subset \{i_{\text{fix}}, i_0, \ldots, i_p\}
$$
The exact same (combinatorial) computation\footnote{
For example, suppose that $i_0 < \ldots < i_p$ is such that
$i_{\text{fix}} \not \in \{i_0, \ldots, i_p\}$ and that
$i_{a - 1} < i_{\text{fix}} < i_a$ for some $1 \leq a \leq p$. Then we have
\begin{align*}
&
(dh + hd)(s)_{i_0 \ldots i_p} \\
& =
\sum\nolimits_{j = 0}^p
(-1)^j
h(s)_{i_0 \ldots \hat i_j \ldots i_p}
+
(-1)^a d(s)_{i_0 \ldots i_{a - 1} i_{\text{fix}} i_a \ldots i_p}\\
& =
\sum\nolimits_{j = 0}^{a - 1}
(-1)^{j + a - 1}
s_{i_0 \ldots \hat i_j \ldots i_{a - 1} i_{\text{fix}} i_a \ldots i_p} +
\sum\nolimits_{j = a}^p
(-1)^{j + a}
s_{i_0 \ldots i_{a - 1} i_{\text{fix}} i_a \ldots \hat i_j \ldots i_p}
+ \\
&
\sum\nolimits_{j = 0}^{a - 1}
(-1)^{a + j}
s_{i_0 \ldots \hat i_j \ldots i_{a - 1} i_{\text{fix}} i_a \ldots i_p} +
(-1)^{2a} s_{i_0 \ldots i_p} +
\sum\nolimits_{j = a}^p
(-1)^{a + j + 1}
s_{i_0 \ldots i_{a - 1} i_{\text{fix}} i_a \ldots \hat i_j \ldots i_p}
\\
& =
s_{i_0 \ldots i_p}
\end{align*}
as desired. The other cases are similar.}
as in the proof of Lemma \ref{lemma-cech-cohomology-quasi-coherent-trivial}
shows that
$$
(hd + dh)(s)_{i_0 \ldots i_p}
=
s_{i_0 \ldots i_p}
$$
Hence we see that the identity map of the complex
$\check{\mathcal{C}}^\bullet(\vec{e})$ is homotopic to zero
which implies that it is acyclic.
\end{proof}

\noindent
In the following lemma we are going to use the pairing of free
$R$-modules
$$
R[T_0, \ldots, T_n]
\times
\frac{1}{T_0 \ldots T_n} R[\frac{1}{T_0}, \ldots, \frac{1}{T_n}]
\longrightarrow
R
$$
which is defined by the rule
$$
(f, g)
\longmapsto
\text{coefficient of }
\frac{1}{T_0 \ldots T_n}
\text{ in }fg.
$$
In other words, the basis element $T_0^{e_0} \ldots T_n^{e_n}$ pairs
with the basis element $T_0^{d_0} \ldots T_n^{d_n}$ to give $1$ if and only
if $e_i + d_i = -1$ for all $i$, and pairs to zero in all other cases.
Using this pairing we get an identification
$$
\left(\frac{1}{T_0 \ldots T_n} R[\frac{1}{T_0}, \ldots, \frac{1}{T_n}]\right)_d
=
\Hom_R((R[T_0, \ldots, T_n])_{-n - 1 - d}, R)
$$
Thus we can reformulate the result of
Lemma \ref{lemma-cohomology-projective-space-over-ring} as saying that
\begin{equation}
\label{equation-identify}
H^q(\mathbf{P}^n, \mathcal{O}_{\mathbf{P}^n_R}(d)) =
\left\{
\begin{matrix}
(R[T_0, \ldots, T_n])_d & \text{if} & q = 0 \\
0 & \text{if} & q \not = 0, n \\
\Hom_R((R[T_0, \ldots, T_n])_{-n - 1 - d}, R)
& \text{if} & q = n
\end{matrix}
\right.
\end{equation}

\begin{lemma}
\label{lemma-identify-functorially}
The identifications of Equation (\ref{equation-identify}) are
compatible with base change w.r.t.\ ring maps $R \to R'$.
Moreover, for any $f \in R[T_0, \ldots, T_n]$ homogeneous
of degree $m$ the map multiplication by $f$
$$
\mathcal{O}_{\mathbf{P}^n_R}(d)
\longrightarrow
\mathcal{O}_{\mathbf{P}^n_R}(d + m)
$$
induces the map on the cohomology group via the identifications
of Equation (\ref{equation-identify}) which is multiplication by
$f$ for $H^0$ and the contragredient of multiplication by $f$
$$
(R[T_0, \ldots, T_n])_{-n - 1 - (d + m)}
\longrightarrow
(R[T_0, \ldots, T_n])_{-n - 1 - d}
$$
on $H^n$.
\end{lemma}

\begin{proof}
Suppose that $R \to R'$ is a ring map.
Let $\mathcal{U}$ be the standard affine open covering of $\mathbf{P}^n_R$,
and let $\mathcal{U}'$ be the standard affine open covering of
$\mathbf{P}^n_{R'}$. Note that $\mathcal{U}'$ is the pullback of the covering
$\mathcal{U}$ under the canonical morphism
$\mathbf{P}^n_{R'} \to \mathbf{P}^n_R$. Hence there
is a map of {\v C}ech complexes
$$
\gamma :
\check{\mathcal{C}}_{ord}^\bullet(\mathcal{U},
\mathcal{O}_{\mathbf{P}_R}(d))
\longrightarrow
\check{\mathcal{C}}_{ord}^\bullet(\mathcal{U}',
\mathcal{O}_{\mathbf{P}_{R'}}(d))
$$
which is compatible with the map on cohomology by
Cohomology, Lemma \ref{cohomology-lemma-functoriality-cech}.
It is clear from the computations in the proof of
Lemma \ref{lemma-cohomology-projective-space-over-ring}
that this map of {\v C}ech complexes is compatible with the identifications
of the cohomology groups in question. (Namely the basis elements for
the {\v C}ech complex over $R$ simply map to the corresponding basis elements
for the {\v C}ech complex over $R'$.) Whence the first statement of the lemma.

\medskip\noindent
Now fix the ring $R$ and consider two homogeneous polynomials
$f, g \in R[T_0, \ldots, T_n]$ both of the same degree $m$.
Since cohomology is an additive functor, it is clear that the
map induced by multiplication by $f + g$ is the same as the sum
of the maps induced by multiplication by $f$ and the map induced
by multiplication by $g$. Moreover, since cohomology is a functor,
a similar result holds for multiplication by a product $fg$ where
$f, g$ are both homogeneous (but not necessarily of the same degree).
Hence to verify the second statement of the lemma it suffices to
prove this when $f = x \in R$ or when $f = T_i$.
In the case of multiplication by an element $x \in R$ the result
follows since every cohomology groups or complex in sight has the
structure of an $R$-module or complex of $R$-modules.
Finally, we consider the case of multiplication by $T_i$
as a $\mathcal{O}_{\mathbf{P}^n_R}$-linear map
$$
\mathcal{O}_{\mathbf{P}^n_R}(d)
\longrightarrow
\mathcal{O}_{\mathbf{P}^n_R}(d + 1)
$$
The statement on $H^0$ is clear. For the statement on $H^n$
consider multiplication by $T_i$ as a map on {\v C}ech complexes
$$
\check{\mathcal{C}}_{ord}^\bullet(\mathcal{U},
\mathcal{O}_{\mathbf{P}_R}(d))
\longrightarrow
\check{\mathcal{C}}_{ord}^\bullet(\mathcal{U},
\mathcal{O}_{\mathbf{P}_R}(d + 1))
$$
We are going to use the notation introduced in the proof of
Lemma \ref{lemma-cohomology-projective-space-over-ring}.
We consider the effect of multiplication by $T_i$
in terms of the decompositions
$$
\check{\mathcal{C}}_{ord}^\bullet(\mathcal{U}, \mathcal{O}_{\mathbf{P}_R}(d))
=
\bigoplus\nolimits_{\vec{e} \in \mathbf{Z}^{n + 1}, \ \sum e_i = d}
\check{\mathcal{C}}^\bullet(\vec{e})
$$
and
$$
\check{\mathcal{C}}_{ord}^\bullet(\mathcal{U},
\mathcal{O}_{\mathbf{P}_R}(d + 1))
=
\bigoplus\nolimits_{\vec{e} \in \mathbf{Z}^{n + 1}, \ \sum e_i = d + 1}
\check{\mathcal{C}}^\bullet(\vec{e})
$$
It is clear that it maps the subcomplex
$\check{\mathcal{C}}^\bullet(\vec{e})$ to the subcomplex
$\check{\mathcal{C}}^\bullet(\vec{e} + \vec{b}_i)$ where
$\vec{b}_i = (0, \ldots, 0, 1, 0, \ldots, 0))$ the $i$th basis vector.
In other words, it maps the summand of $H^n$ corresponding to
$\vec{e}$ with $e_i < 0$ and $\sum e_i = d$
to the summand of $H^n$ corresponding to
$\vec{e} + \vec{b}_i$ (which is zero if $e_i + b_i \geq 0$).
It is easy to see that this corresponds exactly to the action
of the contragredient of multiplication by $T_i$ as a map
$$
(R[T_0, \ldots, T_n])_{-n - 1 - (d + 1)}
\longrightarrow
(R[T_0, \ldots, T_n])_{-n - 1 - d}
$$
This proves the lemma.
\end{proof}

\noindent
Before we state the relative version we need some notation.
Namely, recall that $\mathcal{O}_S[T_0, \ldots, T_n]$ is a graded
$\mathcal{O}_S$-module where each $T_i$ is homogeneous of degree $1$.
Denote $(\mathcal{O}_S[T_0, \ldots, T_n])_d$ the degree $d$ summand.
It is a finite locally free sheaf of rank $\binom{n + d}{d}$ on $S$.

\begin{lemma}
\label{lemma-cohomology-projective-space-over-base}
Let $S$ be a scheme.
Let $n \geq 0$ be an integer.
Consider the structure morphism
$$
f : \mathbf{P}^n_S \longrightarrow S.
$$
We have
$$
R^qf_*(\mathcal{O}_{\mathbf{P}^n_S}(d)) =
\left\{
\begin{matrix}
(\mathcal{O}_S[T_0, \ldots, T_n])_d & \text{if} & q = 0 \\
0 & \text{if} & q \not = 0, n \\
\SheafHom_{\mathcal{O}_S}(
(\mathcal{O}_S[T_0, \ldots, T_n])_{- n - 1 - d}, \mathcal{O}_S)
& \text{if} & q = n
\end{matrix}
\right.
$$
\end{lemma}

\begin{proof}
Omitted. Hint: This follows since the identifications in
(\ref{equation-identify}) are compatible with affine base change
by Lemma \ref{lemma-identify-functorially}.
\end{proof}

\noindent
Next we state the version for projective bundles associated to finite locally
free sheaves. Let $S$ be a scheme. Let $\mathcal{E}$ be a finite locally
free $\mathcal{O}_S$-module of constant rank $n + 1$, see
Modules, Section \ref{modules-section-locally-free}.
In this case we think of $\text{Sym}(\mathcal{E})$ as a graded
$\mathcal{O}_S$-module where $\mathcal{E}$ is the graded part of degree $1$.
And $\text{Sym}^d(\mathcal{E})$ is the degree $d$ summand.
It is a finite locally free sheaf of rank $\binom{n + d}{d}$ on $S$.
Recall that our normalization is that
$$
\pi :
\mathbf{P}(\mathcal{E})
=
\underline{\text{Proj}}_S(\text{Sym}(\mathcal{E}))
\longrightarrow
S
$$
and that there are natural maps
$\text{Sym}^d(\mathcal{E}) \to \pi_*\mathcal{O}_{\mathbf{P}(\mathcal{E})}(d)$.

\begin{lemma}
\label{lemma-cohomology-projective-bundle}
Let $S$ be a scheme. Let $n \geq 1$.
Let $\mathcal{E}$ be a finite locally
free $\mathcal{O}_S$-module of constant rank $n + 1$.
Consider the structure morphism
$$
\pi : \mathbf{P}(\mathcal{E}) \longrightarrow S.
$$
We have
$$
R^q\pi_*(\mathcal{O}_{\mathbf{P}(\mathcal{E})}(d)) =
\left\{
\begin{matrix}
\text{Sym}^d(\mathcal{E}) & \text{if} & q = 0 \\
0 & \text{if} & q \not = 0, n \\
\SheafHom_{\mathcal{O}_S}(
\text{Sym}^{- n - 1 - d}(\mathcal{E})
\otimes_{\mathcal{O}_S}
\wedge^{n + 1}\mathcal{E},
\mathcal{O}_S)
& \text{if} & q = n
\end{matrix}
\right.
$$
These identifications are compatible with base change and
isomorphism between locally free sheaves.
\end{lemma}

\begin{proof}
Consider the canonical map
$$
\pi^*\mathcal{E} \longrightarrow \mathcal{O}_{\mathbf{P}(\mathcal{E})}(1)
$$
and twist down by $1$ to get
$$
\pi^*(\mathcal{E})(-1) \longrightarrow \mathcal{O}_{\mathbf{P}(\mathcal{E})}
$$
This is a surjective map from a locally free rank $n + 1$ sheaf onto
the structure sheaf. Hence the corresponding Koszul complex is
exact (More on Algebra, Lemma
\ref{more-algebra-lemma-homotopy-koszul-abstract}).
In other words there is an exact complex
$$
0 \to
\pi^*(\wedge^{n + 1}\mathcal{E})(-n - 1) \to
\ldots \to
\pi^*(\wedge^i\mathcal{E})(-i) \to
\ldots \to
\pi^*\mathcal{E}(-1) \to
\mathcal{O}_{\mathbf{P}(\mathcal{E})} \to 0
$$
We will think of the term $\pi^*(\wedge^i\mathcal{E})(-i)$ as being
in degree $-i$.
We are going to compute the higher direct images
of this acyclic complex using the first spectral sequence of
Derived Categories, Lemma \ref{derived-lemma-two-ss-complex-functor}.
Namely, we see that there is a spectral sequence with terms
$$
E_1^{p, q} = R^q\pi_*\left(\pi^*(\wedge^{-p}\mathcal{E})(p)\right)
$$
converging to zero! By the projection formula
(Cohomology, Lemma \ref{cohomology-lemma-projection-formula})
we have
$$
E_1^{p, q} = \wedge^{-p} \mathcal{E} \otimes_{\mathcal{O}_S}
R^q\pi_*\left(\mathcal{O}_{\mathbf{P}(\mathcal{E})}(p)\right).
$$
Note that locally on $S$ the sheaf $\mathcal{E}$ is trivial,
i.e., isomorphic to $\mathcal{O}_S^{\oplus n + 1}$, hence locally on
$S$ the morphism $\mathbf{P}(\mathcal{E}) \to S$ can be identified
with $\mathbf{P}^n_S \to S$. Hence
locally on $S$ we can use the result of Lemmas
\ref{lemma-cohomology-projective-space-over-ring},
\ref{lemma-identify-functorially}, or
\ref{lemma-cohomology-projective-space-over-base}.
It follows that $E_1^{p, q} = 0$ unless $(p, q)$ is $(0, 0)$
or $(-n - 1, n)$. The nonzero terms are
\begin{align*}
E_1^{0, 0} & = \pi_*\mathcal{O}_{\mathbf{P}(\mathcal{E})} = \mathcal{O}_S \\
E_1^{-n - 1, n} & =
R^n\pi_*\left(\pi^*(\wedge^{n + 1}\mathcal{E})(-n - 1)\right) =
\wedge^{n + 1}\mathcal{E} \otimes_{\mathcal{O}_S}
R^n\pi_*\left(\mathcal{O}_{\mathbf{P}(\mathcal{E})}(-n - 1)\right)
\end{align*}
Hence there can only be one nonzero
differential in the spectral sequence namely the map
$d_{n + 1}^{-n - 1, n} : E_{n + 1}^{-n - 1, n} \to E_{n + 1}^{0, 0}$
which has to be an isomorphism (because the spectral sequence converges
to the $0$ sheaf). Thus $E_1^{p, q} = E_{n + 1}^{p, q}$ and
we obtain a canonical isomorphism
$$
\wedge^{n + 1}\mathcal{E} \otimes_{\mathcal{O}_S}
R^n\pi_*\left(\mathcal{O}_{\mathbf{P}(\mathcal{E})}(-n - 1)\right) =
R^n\pi_*\left(\pi^*(\wedge^{n + 1}\mathcal{E})(-n - 1)\right)
\xrightarrow{d_{n + 1}^{-n - 1, n}}
\mathcal{O}_S
$$
Since $\wedge^{n + 1}\mathcal{E}$ is an invertible
sheaf, this implies that
$R^n\pi_*\mathcal{O}_{\mathbf{P}(\mathcal{E})}(-n - 1)$ is invertible
as well and canonically isomorphic to the inverse of
$\wedge^{n + 1}\mathcal{E}$. In other words we have proved the case
$d = - n - 1$ of the lemma.

\medskip\noindent
Working locally on $S$ we see immediately from the computation of
cohomology in Lemmas \ref{lemma-cohomology-projective-space-over-ring},
\ref{lemma-identify-functorially}, or
\ref{lemma-cohomology-projective-space-over-base} the statements on
vanishing of the lemma. Moreover the result on $R^0\pi_*$ is clear
as well, since there are canonical maps
$\text{Sym}^d(\mathcal{E}) \to \pi_* \mathcal{O}_{\mathbf{P}(\mathcal{E})}(d)$
for all $d$. It remains to show that the description of
$R^n\pi_*\mathcal{O}_{\mathbf{P}(\mathcal{E})}(d)$ is correct
for $d < -n - 1$. In order to do this we consider the map
$$
\pi^*(\text{Sym}^{-d - n - 1}(\mathcal{E}))
\otimes_{\mathcal{O}_{\mathbf{P}(\mathcal{E})}}
\mathcal{O}_{\mathbf{P}(\mathcal{E})}(d)
\longrightarrow
\mathcal{O}_{\mathbf{P}(\mathcal{E})}(-n - 1)
$$
Applying $R^n\pi_*$ and the projection formula (see above) we get a map
$$
\text{Sym}^{-d - n - 1}(\mathcal{E})
\otimes_{\mathcal{O}_S}
R^n\pi_*(\mathcal{O}_{\mathbf{P}(\mathcal{E})}(d))
\longrightarrow
R^n\pi_*\mathcal{O}_{\mathbf{P}(\mathcal{E})}(-n - 1) =
(\wedge^{n + 1}\mathcal{E})^{\otimes -1}
$$
(the last equality we have shown above).
Again by the local calculations of Lemmas
\ref{lemma-cohomology-projective-space-over-ring},
\ref{lemma-identify-functorially}, or
\ref{lemma-cohomology-projective-space-over-base}
it follows that this map induces a perfect pairing between
$R^n\pi_*(\mathcal{O}_{\mathbf{P}(\mathcal{E})}(d))$ and
$\text{Sym}^{-d - n - 1}(\mathcal{E}) \otimes \wedge^{n + 1}(\mathcal{E})$
as desired.
\end{proof}

















